% BT1408 -- manuscript insert: the frame cross-matching and its cokernel.
%
% Promotion candidate for w33_paper.tex, to sit beside the existing spectral
% material on the 1 + 24 + 15 split.  Every numeric claim below has a GAP
% certificate (data/w33_pass1390_frame_cross_matching.txt,
% data/w33_pass1392_incidence_snf.json, data/w33_pass1397_cokernel_module.txt,
% data/w33_pass1405_1406.txt) and the uniqueness lemma has a Lean module
% (formal/W33/Pass1390FrameCrossMatching.lean).
%
% NOT included on purpose: any E_8 reading of the 240.  Pass 1012 owns the
% obstruction, and the 240 here is the edge set of W(3,3).

% HOST-INDEPENDENT PREAMBLE.  This insert is \input by BOTH root manuscripts and
% they do not define the same theorem environments: w33_paper.tex has
% lemma/proposition/remark, photonic_holonet.tex does not, so promoting this file
% into the Holonet broke that build with "Environment lemma undefined".  Define
% only what the host is missing, so the insert is portable either way.
\makeatletter
\@ifundefined{theorem}{\newtheorem{theorem}{Theorem}}{}
\@ifundefined{lemma}{\newtheorem{lemma}[theorem]{Lemma}}{}
\@ifundefined{proposition}{\newtheorem{proposition}[theorem]{Proposition}}{}
\@ifundefined{remark}{\newtheorem{remark}[theorem]{Remark}}{}
\makeatother
% ...and the operator macros, for the same reason: \PSp and \Aut are defined in
% w33_paper.tex only.
\providecommand{\PSp}{\mathrm{PSp}}
\providecommand{\PGSp}{\mathrm{PGSp}}
\providecommand{\Aut}{\mathrm{Aut}}

\subsection*{The frame cross-matching and the \texorpdfstring{$(-4)$}{(-4)}-eigenspace}

Call a \emph{frame} an unordered pair $\{L,L'\}$ of disjoint totally isotropic
lines of $W(3,3)$; there are $540$ of them, and the stabiliser of a frame in
$\PSp(4,3)$ is $C_2\times S_4\cong O_h$ of order $48$.

\begin{lemma}[Faithful tetrahedral action]\label{lem:frameA4}
Let $H$ be the stabiliser of a frame $\{L,L'\}$ and $A_4=H'$ its derived
subgroup. Then $A_4$ acts faithfully, as the natural degree-four alternating
group, on the four points of $L$ and on the four points of $L'$, with orbits
$[4,4]$ on the eight points of the frame.
\end{lemma}

\begin{proposition}[Canonical cross-matching]\label{prop:crossmatching}
For every frame there is exactly one $A_4$-equivariant bijection $L\to L'$, and
it is invariant under the full stabiliser $H$, not merely under $A_4$.
\end{proposition}

\begin{proof}
A faithful action on both four-element sets embeds $A_4$ into
$\operatorname{Sym}(L)\times\operatorname{Sym}(L')$ with both projections onto,
so its image is the graph of an isomorphism and an equivariant bijection exists.
Uniqueness is general: on a transitive $A$-set an equivariant bijection is
determined by the image of one point, since
$f(a\cdot x_0)=a\cdot f(x_0)$ for all $a$. Invariance under $H$ then follows
because $H$ permutes the set of $A_4$-equivariant bijections, which is a
singleton. Exhaustive verification over all $540$ frames is recorded in the
certificate.
\end{proof}

Write $\mu(\{L,L'\})\subseteq E$ for the resulting set of four \emph{cross-pairs}.

\begin{theorem}[The matching lands on edges, $9$ to $1$]\label{thm:crossedges}
Every cross-pair is an edge of $W(3,3)$; the $540$ matchings use exactly the
$240$ edges, each edge lying in exactly nine of them; and $\mu$ is injective.
\end{theorem}

\begin{proof}
The $240$ cross-pairs form a single $\Aut$-orbit with point stabiliser of order
$108=C_3\times S_3\times S_3$, and $|\PSp(4,3)|/240=108$; uniformity of the
multiplicity follows, and $540\cdot 4/240=9$. Injectivity holds because the
endpoints of the four cross-pairs are exactly $L\cup L'$ and each line is a
maximal clique of the collinearity graph, so the eight points recover the frame.
That every cross-pair is collinear is not forced by the construction --- $L$ and
$L'$ are disjoint --- and is verified.
\end{proof}

Let $M\colon \mathbb{Z}^{540}\to\mathbb{Z}^{240}$ send a frame to the sum of its
four cross-pairs.

\begin{theorem}[Cokernel]\label{thm:crosscoker}
$\operatorname{rank} M=225$, the invariant factors are $1^{195},2^{30}$, and
\[
  \operatorname{coker}(M)\;\cong\;\mathbb{Z}^{15}\oplus(\mathbb{Z}/2)^{30}.
\]
Moreover $\operatorname{coker}(M)\otimes\mathbb{Q}$ is \emph{irreducible} of
degree $15$, and it is the same constituent that occurs in the $40$-point
permutation module. Equivalently,
\[
  \operatorname{coker}(M)\otimes\mathbb{Q}\;\cong\;\ker(A+4I),
\]
the $(-4)$-eigenspace of the adjacency operator.
\end{theorem}

The last identification is the point. The $40$-point module splits as
$1+15+24$, matching the eigenvalue multiplicities of $12,-4,2$; among those three
blocks it is the $(-4)$-block alone that the frame geometry reproduces, and it
does so through an explicitly named equivariant map rather than through an
equality of dimensions.

\begin{remark}[What is \emph{not} claimed]\label{rem:crossscope}
Theorem~\ref{thm:crosscoker} is rational. Integrally the situation differs: $2$
is the only bad prime ($\operatorname{rank}_{\mathbb{F}_2}M=195$ while
$\operatorname{rank}_{\mathbb{F}_3}M=\operatorname{rank}_{\mathbb{F}_5}M=225$),
and the mod-$2$ quotient is \emph{not} irreducible --- its composition factors
have dimensions $1,1,1,6,8,14,14$. So the degree-$15$ rational constituent
reduces to $1+14$ modulo $2$, and the torsion contributes $1,1,6,8,14$; the two
degree-$14$ factors are isomorphic, so the torsion carries a second copy of the
reduction rather than an unrelated module. This is
consistent with the coalescence theorem: $\operatorname{spec}(A)=\{12,2,-4\}$
collapses completely modulo $2$, while modulo $3$ and $5$ it does not.

The corresponding statement for the signed-turn operator is a \emph{theorem},
established in Passes 1416--1420 and not here. With $N$ the unsigned point--edge
incidence and $\partial$ the oriented one,
\[
  F=\tfrac{1}{16}\,\partial^{\mathsf T}(A-12I)(A-2I)N
\]
is integral and satisfies $FM^{\mathsf T}=0$, $(K-10I)F=0$ and
$\operatorname{rank}_{\mathbb{Q}}F=15$, so it induces an equivariant isomorphism
$\operatorname{coker}(M)\otimes\mathbb{Q}\cong\ker(K-10I)$. A dimension-only
argument does \emph{not} suffice, and the reason is worth recording: the
$240$-edge \emph{permutation} module contains two non-isomorphic degree-$15$
constituents, but $K$ acts on the orientation-\emph{signed} edge action, so
$\ker(K-10I)$ is not a submodule of that permutation module at all. The
intertwiner supplies the missing twist. Modulo $2$ the same map has rank $14$ and
square zero, and it separates the two abstractly isomorphic degree-$14$ factors:
it selects the nontrivial reduction of the rational $15$, while the second copy
remains in the torsion-side kernel.
\end{remark}

\subsection*{The Hodge decomposition of the edge module, and what it does not
contain}
\label{sec:hodgeblocks}

The clique complex of $W(3,3)$ is three-dimensional: every totally isotropic line
is a $4$-clique, hence a tetrahedron, giving
$(C_0,C_1,C_2,C_3)=(40,240,160,40)$ with $\chi=-80$ and Betti numbers
$(b_0,b_1,b_2,b_3)=(1,81,0,0)$ --- the complex is homotopy equivalent to a
bouquet of $81$ circles.

\begin{theorem}[Hodge blocks of the signed edge module]\label{thm:hodgeblocks}
Under the orientation-signed action of $\Aut$ on the $240$ edge $1$-chains, the
Hodge decomposition is multiplicity-free and splits into six irreducibles:
\[
\underbrace{15\oplus 24}_{\text{exact, }39}\;\oplus\;
\underbrace{81}_{\text{harmonic}}\;\oplus\;
\underbrace{30\oplus 45\oplus 45}_{\text{coexact, }120}\;=\;240 .
\]
The harmonic block is \emph{irreducible} of degree $81$ and is the Steinberg
module. The exact block $15\oplus24$ is precisely the $40$-point permutation
module modulo its trivial constituent, i.e.\ the $(-4)$- and $(+2)$-eigenspaces
of $A$.
\end{theorem}

In gauge-theoretic language this is the pure-gauge / physical / constraint split,
with the Gauss-law count $39=40-1$ realised representation-theoretically and a
physical sector that is a single irreducible.

\begin{proposition}[The blocks behave oppositely under the outer automorphism]
\label{prop:blockchirality}
$\PSp(4,3)$ has index two in $\PGSp(4,3)\cong W(E_6)$. Under that extension:
\[
15,\;24,\;81 \;\longmapsto\; \text{two extensions each (a sign character)},
\qquad
45\oplus45 \;\longmapsto\; \text{a single irreducible of degree } 90 .
\]
Concretely, $\PGSp(4,3)$ has four degree-$15$, two degree-$24$ and two
degree-$81$ irreducibles, and \emph{no} degree-$45$ irreducible but one of degree
$90$.
\end{proposition}

So the gauge and physical blocks are \emph{split} by the outer element --- each
constituent survives and acquires a sign --- while the two halves of the
constraint block are \emph{exchanged} by it. The chirality of the constraint
sector is of a different kind from that of the physical sector, not merely of a
different magnitude.

\begin{remark}[No Hodge star, and where one must come from]\label{rem:nostar}
A Hodge star requires $C_k\cong C_{3-k}$. Here $C_0=C_3=40$ but
$C_1=240\neq160=C_2$, the discrepancy being exactly $|\chi|$; and on the physical
sector $\star$ would map $H^1\to H^2$ with $b_2=0$. So the $81$ physical modes
admit no dual, there is no $F\wedge\star F$, and no variational dynamics can be
written on this complex. The finite geometry supplies the \emph{kinematics} and
withholds the \emph{dynamics}.

In an almost-commutative geometry $M^4\times F$ the volume form, the star and the
variational principle belong to the continuum factor; the finite factor
contributes an algebra and a module. That division is now explicit: the finite
side supplies $\langle I,A,J\rangle$ and the Steinberg module of
Theorem~\ref{thm:hodgeblocks}, and the star must be supplied by $M^4$. The
absence is established by computation ($b_2=0$), not by omission, and cannot be
repaired internally.
\end{remark}

\begin{proposition}[Resolutions of the edge set]\label{prop:crosscover}
There exist sets of $60$ frames whose cross-matchings partition $E$. No cover is
$\Aut$-invariant, since $\Aut$ is transitive on the $540$ frames and a cover uses
$60$. Cover stabilisers have order $2$, $4$ or $8$ (types $C_2$, $C_4$,
$C_2\times C_2$, $C_4\times C_2$, $D_8$), so a cover is close to rigid --- but the
residual symmetry is \emph{not} uniformly diagonal: a $C_2$-stabilised cover fixes
twelve of its sixty frames (orbit profile $1^{12}2^{24}$), whereas the $C_4$
covers fix none. A compiled census (Pass 1505) exhibits $327$ pairwise disjoint
complete $\PSp(4,3)$-orbits, certifying at least
\[
  228\cdot 12960+75\cdot 6480+9\cdot 6480+9\cdot 3240+6\cdot 3240=3\,547\,800
\]
covers; this is a lower bound and no exhaustive count is claimed. Note that the
orbit census is heavily $C_2$-weighted --- about $83\%$ of covers lie in $C_2$
orbits --- so a small sample drawn by a depth-first search, however its branch
order is permuted, is \emph{not} a uniform sample and its observed proportions
estimate nothing.
\end{proposition}
